\documentclass[11pt,a4paper]{article}
\usepackage[T2A]{fontenc}
\usepackage[utf8]{inputenc}
\usepackage[russian]{babel}
\usepackage[margin=2.3cm]{geometry}
\usepackage{graphicx}
\usepackage{booktabs}
\usepackage{xcolor}
\usepackage{hyperref}
\definecolor{copper}{HTML}{C96F4A}
\hypersetup{colorlinks=true, linkcolor=copper, urlcolor=copper}
\setlength{\parskip}{4pt}
\begin{document}
\begin{center}
{\LARGE\bfseries Исследовательская программа: верификационная платформа (E01–E18, O-J)}\\[8pt]
{\large Монография исследовательской программы · Версия 1.0.0}\\[6pt]
{\color{copper}\small ПРОГРАММА УРОВНЯ · VERIFICATION}\\[4pt]
{\small Исаев Исхак Хамзатович · ORCID 0009-0003-7299-0701}\\[2pt]
{\small AB Cloud Quantum Research · DOI 10.5281/zenodo.21825394}
\end{center}

\tableofcontents
\newpage

\section{Исследовательская программа: верификационная платформа}

\noindent\textit{\color{copper}Программа верификации — независимый судья программы: Julia-платформа без общего кода с ядром и процессором, слепые протоколы, SHA-пины, кросс-языковая репликация. Восемнадцать вердиктов PASS; следующая волна — QEC-память, цифровой двойник, интеграция.}

\subsection{1. Цели программы}

\begin{itemize}
  \item перепроверить ядро и процессор независимой реализацией;
  \item провести слепое чтение нулей ζ двумя движками до T = 1000;
  \item выполнить честный бенчмарк компилятора (класс B);
  \item построить QEC-слой и цифровой двойник машины.
\end{itemize}

\subsection{2. Методология}

Детектор и валидатор разделены; выходы пинуются SHA256 до валидации; пороги предрегистрированы. Кросс-языковая репликация: Julia BigFloat ↔ Python mpmath (0.0 расхождений), numpy-реплики (0 расхождений счётчиков). Независимость реализации — принцип дизайна, а не обстоятельство.

\subsection{3. Исполненные кампании E01–E18}

\begin{table}[htbp]\centering\small
\caption{Таблица 3.3. Кампании верификации}
\begin{tabular}{p{0.21\textwidth} p{0.15\textwidth} p{0.60\textwidth}}
\toprule
\textbf{Блок} & \textbf{Кампании} & \textbf{Ключевые результаты}
\midrule
инфраструктура & E01–E10 & контроли, Haar, масштабирование, реплеи, шумы, код, АЛУ, стоимость, считывание \\
чтение ζ & E11 & слепой ζ-ридер: пины, валидация, ноль утечек \\
оптика/нелокальность & E12–E14 & интерферометр; no-signaling; CHSH \\
классификация & E15 & топологическая абляция \\
движки & E16 & независимый Риман–Зигель: ≤ 3.8e−15 \\
флагман & E17 & слепые нули: 649/649 + 639/639; кросс-движок 0.0; N(1000) = 649.0 \\
бенчмарк & E18 & класс B: ×59 000; 2q KAK вне класса (10 против 3) \\
\bottomrule
\end{tabular}
\end{table}

\subsection{4. Дорожная карта O-J}

\begin{itemize}
  \item O-J1 исполнено: слепое чтение до T = 1000 (E16–E17);
  \item O-J2 исполнено: класс B (E18);
  \item O-J3: QEC-память — опирается на E07 и V21;
  \item O-J4: цифровой двойник процессора;
  \item O-J5: интеграционная кампания; целевые высоты T существенно больше 1000.
\end{itemize}

\subsection{5. Честные границы}

E17 не доказывает RH выше T = 1000, простоту всех нулей и отсутствие нулей вне прямой при γ > 1000. E18 не заявляет превосходства вне ISA-класса. Слепой эрратум γ₁₂ показал, что ошибка может жить в каркасе, а не в детекторе — поэтому эталоны тоже проверяются.

\subsection{6. Разбор ключевых кампаний}

\subsection{6.1. От инфраструктуры к флагманам}

\subsubsection{E11 — слепой ζ-ридер}

Первый полный слепой протокол платформы: детектор без доступа к эталону, пинование SHA256, валидация после фиксации. Кампания отработала саму механику слепоты на живой задаче — чтение нулей меньшей высоты — и показала ноль утечек. Все последующие слепые протоколы (E16, E17) наследуют её схему.

\subsubsection{E16 — независимый Риман–Зигель}

Реализация формулы Римана–Зигеля (версия Ариаса де Рейны) без внешних пакетов: главная сумма плюс интегральная поправка. Согласие с эталоном ≤ 3.8e−15 при пороге 1e−9; пол ошибки 1e−17. Смысл: платформа владеет вторым независимым движком вычисления Z(t), не зависящим от явного метода.

\subsubsection{E17 — слепые нули T=1000}

Флагман платформы: два независимых детектора читают нули до T = 1000 без доступа к эталону; 649/649 (ЭМ) и 639/639 (РС), кросс-движок 0.0 побитово, N(1000) = 649.0 точно, устойчивость к сдвигу сетки — 0. Слепой эрратум γ₁₂ встроен как негативный контроль: детектор опроверг рукотворную константу каркаса. Vердикт честно фиксирует: RH выше T = 1000 не доказывается.

\subsubsection{E18 — класс B}

Честный бенчмарк компилятора: один ISA, один критерий, пять сильных baseline. Итог: ×59 000 по времени при равной точности; 2q-граница 10 против KAK-оптимума 3 — вне класса, заявлено явно. numpy-репликация: 0 расхождений счётчиков. Методология E18 становится стандартом P5 процессорной дорожной карты.

\subsubsection{E12–E14 — оптика и нелокальность}

Блок интерферометрии и нелокальности: интерферометр (E12), отсутствие сигналинга (E13), CHSH (E14). Эти кампании независимо подтверждают превью V24 процессора и дают независимую оценку корреляций распределённых облаков.

\subsection{7. Как воспроизвести}

\begin{verbatim}
julia modules/verifier_multilang/experiments/run_all.jl all   # E01–E18 (~2.5 мин)
julia modules/verifier_multilang/experiments/verdicts.jl      # сводка вердиктов
# пины и результаты: modules/verifier_multilang/results/
\end{verbatim}

\subsection{8. Пины и целостность}

Каждый слепой выход пинован SHA256 до валидации: e17\_detector\_em.csv, e17\_detector\_rs.csv, e11\_detector\_output.csv, эталон e17\_zetazero\_reference.csv. Файлы пинов лежат рядом с данными; валидатор пакета пересчитывает хеши при каждой сборке. Смена пина после валидации невозможна без следа в истории.

\subsection{9. Что читать дальше}

\begin{itemize}
  \item протоколы: experiments/experiment\_01\_15.md;
  \item мастер-вердикты: results/verdicts\_master.json;
  \item кросс-языковая валидация: results/cross\_language\_validation.md;
  \item ограничения и происхождение: PROVENANCE\_AND\_LIMITATIONS.md.
\end{itemize}

\end{document}
